%! Sine and Cosine Function Values — Unit 3, Lesson 3.3 — Exit Ticket
\documentclass[10pt]{article}
\usepackage{precalculus-article}
\usepackage{precalculus-boxes}
\newcommand{\TallMath}[1]{$\displaystyle #1\rule[-1.4em]{0pt}{3.2em}$}

\begin{document}

\pageheader{Unit 3, Lesson 3.3}{Exit Ticket}
\namedateperiod

\vspace{0.16in}

\begin{notesbox}{Sine and Cosine Function Values}

\textbf{1.} Fill in the exact values. Show the reference angle and the quadrant you used.

\vspace{0.1in}
\begin{tabularx}{\linewidth}{@{} X X @{}}
$\cos\!\left(\dfrac{2\pi}{3}\right) = $ \underline{\hspace{2.2cm}} &
$\sin\!\left(\dfrac{5\pi}{6}\right) = $ \underline{\hspace{2.2cm}} \\[1em]
\end{tabularx}

\textit{Reference angle and quadrant for each:}

\vspace{1.1in}

\textbf{2.} A point $P$ lies on a circle of radius 5 centered at the origin, at angle
$\theta = \dfrac{3\pi}{4}$.

State the exact coordinates of $P$. Show your work.

\vspace{1.1in}

$P = \Bigl($ \underline{\hspace{2.2cm}} $,$ \underline{\hspace{2.2cm}} $\Bigr)$

\end{notesbox}

\end{document}
